\chapter{Touch Probing}
\label{ch:probing}

The probe page turns the touch probe into a guided measuring tool: choose a category, pick an
illustrated operation, set the parameters, run it, read the result. The probing math is not
reimplemented -- the Probe Basic G-code routines ship with LiNC and run as they are.

Probing is available on a mill profile with a probe configured. On the Offsets page, \keycap{F8}
moves into the probe tabs.

\shot{img/ui/probe.png}{The probe page: editable DRO and parameters on the left, the routine
matrix in the centre, the probe and results on the right.}{fig:probe}

\section{The tabs}

\begin{center}
\begin{tabular}{ll}
  \toprule
  \textbf{Tab} & \textbf{What it measures} \\
  \midrule
  \state{OUTSIDE CORNERS} & X/Y faces and corners from outside the part \\
  \state{INSIDE CORNERS} & Faces and corners from inside a pocket \\
  \state{BOSS AND POCKET} & Round and rectangular bosses and pockets: width, centre, diameter \\
  \state{RIDGE AND VALLEY} & Ridges, valleys and their centres \\
  \state{EDGE ANGLE} & The angle of an edge, and applying it as a work rotation \\
  \state{CALIBRATE} & Probe calibration: round and square boss and pocket, offsets, reset \\
  \state{PROBE HELP} & The Probe Basic help text, page by page \\
  \bottomrule
\end{tabular}
\end{center}

\keycap{F8} moves to the next tab and \keycap{PgUp}/\keycap{PgDn} cycle them; \keycap{Esc}
returns to Offsets. The tabs are indicators, not buttons, and take no keyboard focus.

\section{The page}

\begin{itemize}
  \item \textbf{Top left:} an editable work-coordinate DRO. \keycap{X}, \keycap{Y} and
        \keycap{Z} open that axis, \keycap{Enter} writes the work position with
        \texttt{G10 L20}, \keycap{Esc} cancels. Work offsets themselves stay on the Offsets
        page; the header names the active system.
  \item \textbf{Below it:} the parameter card, with the probe tool, the fast, slow and
        traverse feed rates, the maximum X/Y and Z distance, the clearances, the extra depth,
        the edge width and the diameter hints. Arithmetic is accepted, as in the offset cells.
  \item \textbf{Centre:} the illustrated routines of the selected tab. \keycap{1} to
        \keycap{9} highlight a routine. \keycap{Enter} runs it.
  \item \textbf{Right:} the probe symbol and the results of the last run.
\end{itemize}

No routine is selected when the page opens or a tab changes. A number key is the only way to
select one; \keycap{Enter} before that reports that a routine must be selected.

\begin{tipbox}
  The keys that matter are listed in the pane under the navigation column: \state{PROBE KEYS}
  and \state{RESULT KEYS}. After an hour away from the machine, that pane is faster than this
  chapter.
\end{tipbox}

\section{Position and direct: the W key}

\keycap{W} switches what a run does with its result, and the screen says which mode is active:

\begin{description}[font=\normalfont\sffamily\bfseries, leftmargin=3.4cm, style=nextline]
  \item[\state{POSITION}] The default. The routine measures and records. Nothing is written to
        a work offset; \keycap{A} applies the last result when you decide to.
  \item[\state{DIRECT WCO}] The routine writes the active work offset itself, in the issued
        block. It needs a valid active offset, and it is refused when protect is on.
\end{description}

\section{Running a routine}

\keycap{Enter} builds the selected \texttt{o<probe\_...> call} and its argument values and
issues it immediately. The controller is switched to MDI for the block and the previous task
mode is restored when the interpreter returns idle. No physical cycle start is involved.

A run requires:

\begin{itemize}
  \item machine power on and every axis homed,
  \item the probe tool in the spindle,
  \item an idle interpreter and no active program,
  \item valid parameters.
\end{itemize}

A validation error issues nothing and marks the failing condition.

\begin{cautionbox}
  A probe cycle moves the machine. Check the clearance values and the probe tool before
  \keycap{Enter}. The fast feed rate is the approach; the slow feed rate is the contact.
\end{cautionbox}

\section{The results}

Results come from the subroutine outputs and are shown as named slots: \state{X WIDTH},
\state{X CENTRE}, \state{Y WIDTH}, \state{Y CENTRE}, \state{X- FACE}, \state{X+ FACE},
\state{Y- FACE}, \state{Y+ FACE}, \state{Z SURFACE}, \state{AVG DIAM} and
\state{EDGE DELTA}.

\begin{longtable}{L{4.7cm}L{8.1cm}}
  \toprule
  \textbf{Keys} & \textbf{Action} \\
  \midrule
  \endhead
  \keycap{A} & Apply the last result to the active work offset. Protect must be off. \\
  \keycap{U} & Update the probe parameters from the last result. \\
  \chordtwo{Ctrl}{R} & Reset all probe data. \\
  \chordtwo{Ctrl}{X} / \chordtwo{Ctrl}{Y} & Reset the X or the Y data. \\
  \bottomrule
\end{longtable}

The results panel carries no buttons: those keys are the whole surface.

\section{Edge angle and work rotation}

On \state{EDGE ANGLE}, \keycap{R} applies the measured angle as the active work coordinate
rotation. Check the Offsets page afterwards: the rotation appears there, and the preview shows
the part rotated around its zero.

\section{Calibration}

The \state{CALIBRATE} tab measures the probe against a known feature and stores the calibration
offset. Its fields -- known width or diameter, calibration axis and the calibration offset
itself -- live on the same panel.

\begin{itemize}
  \item \keycap{1} to \keycap{4} highlight a calibration routine; \keycap{Enter} runs it.
  \item \keycap{V} switches the calibration between the averaged XY result and the single X or
        Y result.
  \item \chordtwo{Ctrl}{P} reaches the calibration dimensions.
  \item \chordtwo{Ctrl}{R} resets the calibration through the Probe Basic reset routine.
\end{itemize}

\section{When the probe is hit}

If the probe input rises during a move that is not a probe, the controller aborts the move and
reports an operator error. LiNC latches it: a red banner carries the controller's text on every
page, motion and MDI are blocked, and the machine is switched off unless it was paused with
nothing moving.

\keycap{Esc} clears the latched trip from any page. Re-powering afterwards needs a fresh
\pbtn{ON} at the panel. The full behaviour is in chapter~\ref{ch:power}.

\begin{notebox}
  The probe input is also shown in the bottom bar, but only while the probe tool is in the
  spindle: \state{PROBE READY} or \state{PROBE TRIPPED}. There is no probe field for a cutter in
  the spindle, because there is no probe to read.
\end{notebox}
